% Shared preamble for Week NN: Counting-out circles (final).
% Each band file defines \bandlabel and \packetid before % Shared preamble for Week NN: Counting-out circles (final).
% Each band file defines \bandlabel and \packetid before % Shared preamble for Week NN: Counting-out circles (final).
% Each band file defines \bandlabel and \packetid before % Shared preamble for Week NN: Counting-out circles (final).
% Each band file defines \bandlabel and \packetid before \input{common}.
\usepackage[T1]{fontenc}
\usepackage[default]{sourcesanspro}
\usepackage[letterpaper,left=1.5cm,right=1.5cm,top=2.1cm,bottom=2.0cm,
  headheight=16pt,headsep=0.6cm,footskip=0.9cm]{geometry}
\usepackage{fancyhdr}
\usepackage{tikz}
\usetikzlibrary{calc}
\usepackage{xcolor}

\pagestyle{fancy}
\fancyhf{}
\renewcommand{\headrulewidth}{0pt}
\fancyhead[L]{\normalsize Week NN / Counting-out circles / \bandlabel}
\fancyfoot[L]{\small Bellingham Math Circle / Week NN / \packetid}
\fancyfoot[R]{\small\thepage}

\setlength{\parindent}{0pt}
\setlength{\parskip}{0.5\baselineskip}
\raggedright
\hyphenpenalty=10000
\exhyphenpenalty=10000

\definecolor{ringline}{gray}{0.72}
\definecolor{boxline}{gray}{0.55}

\newcommand{\prob}[1]{\textbf{Problem #1:}}

% \ringat{x}{y}{n}{R}{spot radius}{label gap}{label font}: draws inside the
% current tikzpicture. Spots are numbered 1..n clockwise with 1 at the top;
% the numbers sit just outside the spots so counters do not cover them.
\newcommand{\ringat}[7]{%
  \begin{scope}[shift={(#1,#2)}]
    \ifnum#3>1
      \draw[ringline, line width=0.8pt] (0,0) circle (#4);
    \fi
    \foreach \i in {1,...,#3} {
      \pgfmathsetmacro\ang{90-360*(\i-1)/#3}
      \filldraw[fill=white, draw=black, line width=1pt] (\ang:#4) circle (#5);
      \node[font=#7, inner sep=0pt] at (\ang:{#4+#5+#6}) {\i};
    }
  \end{scope}}

% Stand-alone ring picture: \ring{n}{R}{spot radius}{label gap}{label font}.
\newcommand{\ring}[5]{%
  \begin{tikzpicture}[baseline=(current bounding box.center)]
    \ringat{0}{0}{#1}{#2}{#3}{#4}{#5}
  \end{tikzpicture}}

% Working ring for counters (spots 2.6 cm across). Arg: n, R.
\newcommand{\bigring}[2]{\ring{#1}{#2}{1.3}{0.42}{\large\bfseries}}
\newcommand{\bigringat}[4]{\ringat{#1}{#2}{#3}{#4}{1.3}{0.42}{\large\bfseries}}

% A record table. #1 = first-column width (cm), #2 = cell width (cm),
% #3 = cell height (cm), #4 = comma list of column headings,
% #5 = row labels as {label1,label2,...} (first label goes with the headings row).
\newcommand{\rectable}[5]{%
  \begin{tikzpicture}[x=1cm,y=1cm]
    \def\fw{#1}\def\cw{#2}\def\ch{#3}
    \foreach \lab [count=\r from 0] in {#5} {
      \draw[line width=0.7pt] (0,-\r*\ch) rectangle (\fw,-\r*\ch-\ch);
      \node[anchor=west, inner sep=3pt, text width=\fw cm - 6pt, align=left,
            text height=1.6ex, text depth=0.4ex,
            font=\small] at (0,-\r*\ch-\ch/2) {\lab};
      \foreach \h [count=\c from 0] in {#4} {
        \draw[line width=0.7pt] (\fw+\c*\cw,-\r*\ch) rectangle (\fw+\c*\cw+\cw,-\r*\ch-\ch);
        \ifnum\r=0
          \node[font=\bfseries, text height=1.6ex, text depth=0.3ex] at (\fw+\c*\cw+\cw/2,-\ch/2) {\h};
        \fi
      }
    }
  \end{tikzpicture}}

% Answer box: #1 = width (cm), #2 = height (cm), #3 = optional label at top left.
\newcommand{\answerbox}[3]{%
  \begin{tikzpicture}
    \draw[boxline, line width=0.8pt, rounded corners=4pt] (0,0) rectangle (#1,-#2);
    \node[anchor=north west, inner sep=5pt, font=\small] at (0,0) {#3};
  \end{tikzpicture}}

% Full-page mat: a 16-spot ring with 2.6 cm spots (for 1-inch counters),
% about 18.5 cm across with its numbers. Used as a page of its own.
\newcommand{\matsixteen}{%
  \begin{tikzpicture}
    \ringat{0}{0}{16}{7.2}{1.3}{0.42}{\large\bfseries}
  \end{tikzpicture}}
.
\usepackage[T1]{fontenc}
\usepackage[default]{sourcesanspro}
\usepackage[letterpaper,left=1.5cm,right=1.5cm,top=2.1cm,bottom=2.0cm,
  headheight=16pt,headsep=0.6cm,footskip=0.9cm]{geometry}
\usepackage{fancyhdr}
\usepackage{tikz}
\usetikzlibrary{calc}
\usepackage{xcolor}

\pagestyle{fancy}
\fancyhf{}
\renewcommand{\headrulewidth}{0pt}
\fancyhead[L]{\normalsize Week NN / Counting-out circles / \bandlabel}
\fancyfoot[L]{\small Bellingham Math Circle / Week NN / \packetid}
\fancyfoot[R]{\small\thepage}

\setlength{\parindent}{0pt}
\setlength{\parskip}{0.5\baselineskip}
\raggedright
\hyphenpenalty=10000
\exhyphenpenalty=10000

\definecolor{ringline}{gray}{0.72}
\definecolor{boxline}{gray}{0.55}

\newcommand{\prob}[1]{\textbf{Problem #1:}}

% \ringat{x}{y}{n}{R}{spot radius}{label gap}{label font}: draws inside the
% current tikzpicture. Spots are numbered 1..n clockwise with 1 at the top;
% the numbers sit just outside the spots so counters do not cover them.
\newcommand{\ringat}[7]{%
  \begin{scope}[shift={(#1,#2)}]
    \ifnum#3>1
      \draw[ringline, line width=0.8pt] (0,0) circle (#4);
    \fi
    \foreach \i in {1,...,#3} {
      \pgfmathsetmacro\ang{90-360*(\i-1)/#3}
      \filldraw[fill=white, draw=black, line width=1pt] (\ang:#4) circle (#5);
      \node[font=#7, inner sep=0pt] at (\ang:{#4+#5+#6}) {\i};
    }
  \end{scope}}

% Stand-alone ring picture: \ring{n}{R}{spot radius}{label gap}{label font}.
\newcommand{\ring}[5]{%
  \begin{tikzpicture}[baseline=(current bounding box.center)]
    \ringat{0}{0}{#1}{#2}{#3}{#4}{#5}
  \end{tikzpicture}}

% Working ring for counters (spots 2.6 cm across). Arg: n, R.
\newcommand{\bigring}[2]{\ring{#1}{#2}{1.3}{0.42}{\large\bfseries}}
\newcommand{\bigringat}[4]{\ringat{#1}{#2}{#3}{#4}{1.3}{0.42}{\large\bfseries}}

% A record table. #1 = first-column width (cm), #2 = cell width (cm),
% #3 = cell height (cm), #4 = comma list of column headings,
% #5 = row labels as {label1,label2,...} (first label goes with the headings row).
\newcommand{\rectable}[5]{%
  \begin{tikzpicture}[x=1cm,y=1cm]
    \def\fw{#1}\def\cw{#2}\def\ch{#3}
    \foreach \lab [count=\r from 0] in {#5} {
      \draw[line width=0.7pt] (0,-\r*\ch) rectangle (\fw,-\r*\ch-\ch);
      \node[anchor=west, inner sep=3pt, text width=\fw cm - 6pt, align=left,
            text height=1.6ex, text depth=0.4ex,
            font=\small] at (0,-\r*\ch-\ch/2) {\lab};
      \foreach \h [count=\c from 0] in {#4} {
        \draw[line width=0.7pt] (\fw+\c*\cw,-\r*\ch) rectangle (\fw+\c*\cw+\cw,-\r*\ch-\ch);
        \ifnum\r=0
          \node[font=\bfseries, text height=1.6ex, text depth=0.3ex] at (\fw+\c*\cw+\cw/2,-\ch/2) {\h};
        \fi
      }
    }
  \end{tikzpicture}}

% Answer box: #1 = width (cm), #2 = height (cm), #3 = optional label at top left.
\newcommand{\answerbox}[3]{%
  \begin{tikzpicture}
    \draw[boxline, line width=0.8pt, rounded corners=4pt] (0,0) rectangle (#1,-#2);
    \node[anchor=north west, inner sep=5pt, font=\small] at (0,0) {#3};
  \end{tikzpicture}}

% Full-page mat: a 16-spot ring with 2.6 cm spots (for 1-inch counters),
% about 18.5 cm across with its numbers. Used as a page of its own.
\newcommand{\matsixteen}{%
  \begin{tikzpicture}
    \ringat{0}{0}{16}{7.2}{1.3}{0.42}{\large\bfseries}
  \end{tikzpicture}}
.
\usepackage[T1]{fontenc}
\usepackage[default]{sourcesanspro}
\usepackage[letterpaper,left=1.5cm,right=1.5cm,top=2.1cm,bottom=2.0cm,
  headheight=16pt,headsep=0.6cm,footskip=0.9cm]{geometry}
\usepackage{fancyhdr}
\usepackage{tikz}
\usetikzlibrary{calc}
\usepackage{xcolor}

\pagestyle{fancy}
\fancyhf{}
\renewcommand{\headrulewidth}{0pt}
\fancyhead[L]{\normalsize Week NN / Counting-out circles / \bandlabel}
\fancyfoot[L]{\small Bellingham Math Circle / Week NN / \packetid}
\fancyfoot[R]{\small\thepage}

\setlength{\parindent}{0pt}
\setlength{\parskip}{0.5\baselineskip}
\raggedright
\hyphenpenalty=10000
\exhyphenpenalty=10000

\definecolor{ringline}{gray}{0.72}
\definecolor{boxline}{gray}{0.55}

\newcommand{\prob}[1]{\textbf{Problem #1:}}

% \ringat{x}{y}{n}{R}{spot radius}{label gap}{label font}: draws inside the
% current tikzpicture. Spots are numbered 1..n clockwise with 1 at the top;
% the numbers sit just outside the spots so counters do not cover them.
\newcommand{\ringat}[7]{%
  \begin{scope}[shift={(#1,#2)}]
    \ifnum#3>1
      \draw[ringline, line width=0.8pt] (0,0) circle (#4);
    \fi
    \foreach \i in {1,...,#3} {
      \pgfmathsetmacro\ang{90-360*(\i-1)/#3}
      \filldraw[fill=white, draw=black, line width=1pt] (\ang:#4) circle (#5);
      \node[font=#7, inner sep=0pt] at (\ang:{#4+#5+#6}) {\i};
    }
  \end{scope}}

% Stand-alone ring picture: \ring{n}{R}{spot radius}{label gap}{label font}.
\newcommand{\ring}[5]{%
  \begin{tikzpicture}[baseline=(current bounding box.center)]
    \ringat{0}{0}{#1}{#2}{#3}{#4}{#5}
  \end{tikzpicture}}

% Working ring for counters (spots 2.6 cm across). Arg: n, R.
\newcommand{\bigring}[2]{\ring{#1}{#2}{1.3}{0.42}{\large\bfseries}}
\newcommand{\bigringat}[4]{\ringat{#1}{#2}{#3}{#4}{1.3}{0.42}{\large\bfseries}}

% A record table. #1 = first-column width (cm), #2 = cell width (cm),
% #3 = cell height (cm), #4 = comma list of column headings,
% #5 = row labels as {label1,label2,...} (first label goes with the headings row).
\newcommand{\rectable}[5]{%
  \begin{tikzpicture}[x=1cm,y=1cm]
    \def\fw{#1}\def\cw{#2}\def\ch{#3}
    \foreach \lab [count=\r from 0] in {#5} {
      \draw[line width=0.7pt] (0,-\r*\ch) rectangle (\fw,-\r*\ch-\ch);
      \node[anchor=west, inner sep=3pt, text width=\fw cm - 6pt, align=left,
            text height=1.6ex, text depth=0.4ex,
            font=\small] at (0,-\r*\ch-\ch/2) {\lab};
      \foreach \h [count=\c from 0] in {#4} {
        \draw[line width=0.7pt] (\fw+\c*\cw,-\r*\ch) rectangle (\fw+\c*\cw+\cw,-\r*\ch-\ch);
        \ifnum\r=0
          \node[font=\bfseries, text height=1.6ex, text depth=0.3ex] at (\fw+\c*\cw+\cw/2,-\ch/2) {\h};
        \fi
      }
    }
  \end{tikzpicture}}

% Answer box: #1 = width (cm), #2 = height (cm), #3 = optional label at top left.
\newcommand{\answerbox}[3]{%
  \begin{tikzpicture}
    \draw[boxline, line width=0.8pt, rounded corners=4pt] (0,0) rectangle (#1,-#2);
    \node[anchor=north west, inner sep=5pt, font=\small] at (0,0) {#3};
  \end{tikzpicture}}

% Full-page mat: a 16-spot ring with 2.6 cm spots (for 1-inch counters),
% about 18.5 cm across with its numbers. Used as a page of its own.
\newcommand{\matsixteen}{%
  \begin{tikzpicture}
    \ringat{0}{0}{16}{7.2}{1.3}{0.42}{\large\bfseries}
  \end{tikzpicture}}
.
\usepackage[T1]{fontenc}
\usepackage[default]{sourcesanspro}
\usepackage[letterpaper,left=1.5cm,right=1.5cm,top=2.1cm,bottom=2.0cm,
  headheight=16pt,headsep=0.6cm,footskip=0.9cm]{geometry}
\usepackage{fancyhdr}
\usepackage{tikz}
\usetikzlibrary{calc}
\usepackage{xcolor}

\pagestyle{fancy}
\fancyhf{}
\renewcommand{\headrulewidth}{0pt}
\fancyhead[L]{\normalsize Week NN / Counting-out circles / \bandlabel}
\fancyfoot[L]{\small Bellingham Math Circle / Week NN / \packetid}
\fancyfoot[R]{\small\thepage}

\setlength{\parindent}{0pt}
\setlength{\parskip}{0.5\baselineskip}
\raggedright
\hyphenpenalty=10000
\exhyphenpenalty=10000

\definecolor{ringline}{gray}{0.72}
\definecolor{boxline}{gray}{0.55}

\newcommand{\prob}[1]{\textbf{Problem #1:}}

% \ringat{x}{y}{n}{R}{spot radius}{label gap}{label font}: draws inside the
% current tikzpicture. Spots are numbered 1..n clockwise with 1 at the top;
% the numbers sit just outside the spots so counters do not cover them.
\newcommand{\ringat}[7]{%
  \begin{scope}[shift={(#1,#2)}]
    \ifnum#3>1
      \draw[ringline, line width=0.8pt] (0,0) circle (#4);
    \fi
    \foreach \i in {1,...,#3} {
      \pgfmathsetmacro\ang{90-360*(\i-1)/#3}
      \filldraw[fill=white, draw=black, line width=1pt] (\ang:#4) circle (#5);
      \node[font=#7, inner sep=0pt] at (\ang:{#4+#5+#6}) {\i};
    }
  \end{scope}}

% Stand-alone ring picture: \ring{n}{R}{spot radius}{label gap}{label font}.
\newcommand{\ring}[5]{%
  \begin{tikzpicture}[baseline=(current bounding box.center)]
    \ringat{0}{0}{#1}{#2}{#3}{#4}{#5}
  \end{tikzpicture}}

% Working ring for counters (spots 2.6 cm across). Arg: n, R.
\newcommand{\bigring}[2]{\ring{#1}{#2}{1.3}{0.42}{\large\bfseries}}
\newcommand{\bigringat}[4]{\ringat{#1}{#2}{#3}{#4}{1.3}{0.42}{\large\bfseries}}

% A record table. #1 = first-column width (cm), #2 = cell width (cm),
% #3 = cell height (cm), #4 = comma list of column headings,
% #5 = row labels as {label1,label2,...} (first label goes with the headings row).
\newcommand{\rectable}[5]{%
  \begin{tikzpicture}[x=1cm,y=1cm]
    \def\fw{#1}\def\cw{#2}\def\ch{#3}
    \foreach \lab [count=\r from 0] in {#5} {
      \draw[line width=0.7pt] (0,-\r*\ch) rectangle (\fw,-\r*\ch-\ch);
      \node[anchor=west, inner sep=3pt, text width=\fw cm - 6pt, align=left,
            text height=1.6ex, text depth=0.4ex,
            font=\small] at (0,-\r*\ch-\ch/2) {\lab};
      \foreach \h [count=\c from 0] in {#4} {
        \draw[line width=0.7pt] (\fw+\c*\cw,-\r*\ch) rectangle (\fw+\c*\cw+\cw,-\r*\ch-\ch);
        \ifnum\r=0
          \node[font=\bfseries, text height=1.6ex, text depth=0.3ex] at (\fw+\c*\cw+\cw/2,-\ch/2) {\h};
        \fi
      }
    }
  \end{tikzpicture}}

% Answer box: #1 = width (cm), #2 = height (cm), #3 = optional label at top left.
\newcommand{\answerbox}[3]{%
  \begin{tikzpicture}
    \draw[boxline, line width=0.8pt, rounded corners=4pt] (0,0) rectangle (#1,-#2);
    \node[anchor=north west, inner sep=5pt, font=\small] at (0,0) {#3};
  \end{tikzpicture}}

% Full-page mat: a 16-spot ring with 2.6 cm spots (for 1-inch counters),
% about 18.5 cm across with its numbers. Used as a page of its own.
\newcommand{\matsixteen}{%
  \begin{tikzpicture}
    \ringat{0}{0}{16}{7.2}{1.3}{0.42}{\large\bfseries}
  \end{tikzpicture}}
